% !TeX root = ../../Book4.tex
% 纲要标题：### 7.3 拟同构
% 纲要标签：b4:sec:0603
% 纲要位置：计划纲要.md，第 3258 行。

\section{拟同构}
\label{b4:sec:0603}

同调上的可逆性弱于在复形层面构造同伦逆。映射锥提供拟同构的判据，而非分裂正合列说明这两种可逆性为何会分离。

% 纲要内容（仅作写作计划，尚非教材正文）：
%
% * 拟同构的定义
% * 与同伦等价的区别
% * 非分裂正合列产生的基本例子
%

\subsection{拟同构的定义}
\label{b4:subsec:0603:01}

\begin{definition}\label{b4:def:quasi-isomorphism}
设 \(\mathcal A\) 是交换范畴，\(f:C\to D\) 是复形映射。若每个整数 \(n\) 的 \(H^n(f):H^n(C)\to H^n(D)\) 都是同构，则称 \(f\) 为\term[ni-tonggou]{拟同构}[quasi-isomorphism]。
\end{definition}

拟同构必须是已经给定的复形映射。“两个复形的同调对象逐次同构”只给出对象之间的比较，未必指定了能够实现这些同构的复形映射。更不能由此断言原来的某个映射就是拟同构。

\begin{proposition}\label{b4:prop:cone-quasi-isomorphism}
设 \(\mathcal A\) 是交换范畴，\(f:C\to D\) 是 \(\mathcal A\) 中上链复形之间的映射。则 \(f\) 是拟同构，当且仅当 \(\operatorname{Cone}(f)\) 零调。对 \(\mathcal A\) 中任意可复合的复形映射 \(f:C\to D\)、\(g:D\to E\)，若 \(f,g,gf\) 的任意两个为拟同构，则第三个也是拟同构。
\end{proposition}
\begin{proof}
锥的长正合列中有
\[
H^n(C)\xrightarrow{H^n(f)}H^n(D)\to
H^n\operatorname{Cone}(f)\to H^{n+1}(C)
\xrightarrow{H^{n+1}(f)}H^{n+1}(D).
\]
若左右两个 \(H(f)\) 都是同构，中间对象到右侧的箭头既单又为零，故中间对象为零。若所有锥同调均为零，则相关正合段给出 \(H^n(f)\) 既单又满，故为同构。最后，对 \(H^n(f),H^n(g)\) 使用同构的复合与消去性质，再让 \(n\) 遍历全部整数即可。
\end{proof}

\begin{proposition}\label{b4:prop:short-exact-quasi-comparison}
设交换范畴中两条复形短正合列之间给定态射。如果三个竖直复形映射中有两个为拟同构，则第三个也为拟同构。
\end{proposition}
\begin{proof}
由\cref{b4:thm:cohomology-long-exact}得到两条长正合列及其间交换图。对待证映射 \(H^n\) 所在的项，取其前后各两项组成五项正合图。其余四个竖直态射均是另两个复形在某个次数的同调态射，所以都是同构。\cref{b4:thm:five-lemma}证明中间态射也是同构。
\end{proof}

\subsection{与同伦等价的区别}
\label{b4:subsec:0603:02}

同伦等价必为拟同构，因为同伦逆在同调上成为真正的逆；这是\cref{b4:prop:homotopy-cohomology}的直接应用。拟同构只要求同调上可逆，同伦等价还要求构造反向的复形映射及两个同伦，故要求更强。

\begin{proposition}\label{b4:prop:field-complex-splitting}
设 \(k\) 是域，\(C\) 是任意 \(k\)-向量空间复形。把各个 \(H^n(C)\) 组成微分全零的复形 \(H(C)\)，则 \(C\) 与 \(H(C)\) 同伦等价。因此 \(k\)-向量空间复形之间的每个拟同构都是同伦等价。所用同伦等价一般依赖补空间的选择。
\end{proposition}
\begin{proof}
在每个次数选择补空间
\[
Z^n(C)=B^n(C)\oplus V^n,\qquad C^n=Z^n(C)\oplus W^n.
\]
微分的限制
\[
\mathrm{d}_C^n|_{W^n}:W^n\xrightarrow{\sim}B^{n+1}(C)
\]
是同构，因为其核与 \(W^n\) 的交为零且其像就是全部边界。令 \(q^n:Z^n(C)\to H^n(C)\) 为商映射，则其限制 \(\rho^n=q^n|_{V^n}\) 也是同构。\(H^n(C)\) 本身是商空间；通过 \((\rho^n)^{-1}\)，才能将它识别为选定的代表子空间 \(V^n\subseteq C^n\)。

利用分解 \(C^n=B^n(C)\oplus V^n\oplus W^n\)，定义
\[
i^n:H^n(C)\longrightarrow C^n,\qquad
i^n(\eta)=(\rho^n)^{-1}(\eta),
\]
以及
\[
p^n:C^n\longrightarrow H^n(C),\qquad
p^n(b+v+w)=\rho^n(v).
\]
\(i^n\) 的像在余圈中，而 \(p^{n+1}\) 在微分的像 \(B^{n+1}(C)\) 上为零，故它们组成复形映射 \(i:H(C)\to C\)、\(p:C\to H(C)\)，且 \(pi=1_{H(C)}\)。再定义
\[
h^n(b+v+w)
=\bigl(\mathrm{d}_C^{n-1}|_{W^{n-1}}\bigr)^{-1}(b).
\]
逐个检查这三个分量，得到
\[
\mathrm{d}_C^{n-1}h^n+h^{n+1}\mathrm{d}_C^n
=1_{C^n}-i^np^n.
\]
于是 \(ip\simeq1_C\)，这给出所需同伦等价。

对拟同构 \(f:C\to D\)，给两者分别作上述选择。\(i_C\) 选取同调类的余圈代表，而 \(p_D\) 在余圈上的限制就是取同调类，因此 \(p_Dfi_C=H(f)\) 是微分为零复形之间的同构。取 \(g=i_C H(f)^{-1}p_D\)。由同伦对复合的相容性，
\[
gf\simeq gf i_Cp_C=i_Cp_C\simeq1_C,
\qquad
fg\simeq i_Dp_Dfg=i_Dp_D\simeq1_D.
\]
所以 \(g\) 是 \(f\) 的同伦逆。选择补空间改变了 \(i,p,h\)，所以这不是无须选择的复形分解。
\end{proof}

上述证明使用的是向量空间每个短正合列都分裂。对一般模，核和像可能没有补模；不能只把 \(k\) 改成任意环便保留结论。

\subsection{非分裂正合列产生的基本例子}
\label{b4:subsec:0603:03}

\begin{example}\label{b4:exm:resolution-quasi-not-homotopy}
令 \(P\) 为次数 \(-1,0\) 的复形 \(\mathbb Z\xrightarrow{2}\mathbb Z\)，令 \(M=(\mathbb Z/2)[0]\)。商映射给出增广 \(\varepsilon:P\to M\)，且 \(H^{-1}(P)=0\)、\(H^0(P)=\mathbb Z/2\)，所以它是拟同构。

\ProseParagraphBreak
但任意复形映射 \(g:M\to P\) 的零次分量 \(\mathbb Z/2\to\mathbb Z\) 都为零，故 \(g=0\)。在集中于零次的复形 \(M\) 上，所有可能的降次同伦也为零，因此 \(\varepsilon g=0\) 不可能同伦于 \(1_M\)。所以 \(\varepsilon\) 不是同伦等价。
\end{example}

\begin{example}\label{b4:exm:additive-not-quasi-preserving}
对上一例逐次应用加性函子 \(\operatorname{Hom}_{\mathbb Z}(\mathbb Z/2,-)\)。由于 \(\operatorname{Hom}_{\mathbb Z}(\mathbb Z/2,\mathbb Z)=0\)，所得 \(F(P)\) 是零复形；而 \(F(M)\) 集中于零次且等于 \(\mathbb Z/2\)。因此 \(F(\varepsilon)\) 不是拟同构。
\end{example}

加性函子会保持已经给定的链同伦，因为它保持同伦公式中的加法与复合，却未必保持拟同构。后面先取合适消解再施加函子，正是为了处理这种差别。对非分裂短正合列形成的零调复形，零复形到它的映射也是拟同构；若它不是可缩复形，就同样不会成为同伦等价。
